\subsection{Winding Number}

we know what winding number is. note that:
\begin{enumerate}[1)]
	\item $\gamma(t)=(\cos(2\pi nt),\sin(2\pi nt))$ has winding num $n$
	\item winding num is inv under deformations that stay in the punctured plane.
\end{enumerate}
number 2 suggests we can define $\int_{\gamma}$(closed form).

\begin{defn}
	the \textbf{winding number} of a loop $\gamma$ is given by:
	\begin{equation*}
		w(\gamma) \ = \ \int_{S^{1}}\gamma^{*}\left(\frac{\d\theta}{2\pi}\right)
	\end{equation*}
\end{defn}
note $\theta$ is as in polar coords, but $\theta$ is not a smooth fn
$\bR^{2}\sm\set{0}\gto\bR$, so $\d\theta$ is not exact.

compute: if $x=r\cos\theta,r^{2}=x^{2}+y^{2}$, then:
\begin{equation*}
	\d x \ = \ -r\sin\theta\d\theta+\d r\cos\theta \qquad
	2r\d r \ = \ 2x\d x+2y\d y
\end{equation*}
therefore:
\begin{equation*}
	\d\theta \ = \ \frac{x\d y-y\d x}{x^{2}+y^{2}}
\end{equation*}
check its closed.
this satisfies 1) by computation, and 2) by properties of closed forms.

moreover:
\begin{enumerate}[1)]
	\item $w(\gamma)\in\bZ$
	\item two loops can be deformed into each other iff they have the same winding num
	\item joining loops by a path results in a loop whose winding num is the
		sum of the others' winding nums
\end{enumerate}

\begin{thm}[title=Poincare's Theorem]
	let $\Sigma$ be cpt oriented surface w $\p$, $X$ VF on $\Sigma$ w \textit{isolated zeroes}
	(or equivalently by cptness, finitely many). then:
	\begin{equation*}
		\sum_{\trm{0's}}\trm{index}_{p}(X) - \sum_{C \trm{ bdry circles}}\trm{rot}_{C}(X) \ = \ \chi(\Sigma)
	\end{equation*}
	where $\chi$ is the euler characteristic.
\end{thm}
heuristic: $\fa X$ VF, sth abt how it behaves near 0's + sth abt how it behaves
at $\p\Sigma$ = topological invariant.

the euler characteristic has many defns. for instance:
\begin{equation*}
	\chi(M) \ = \ \sum_{i=0}^{m}(-1)^{i}\trm{dim}(H^{i}(M))
\end{equation*}
we take an axiomatic route.

\begin{defn}
	the euler characteristic is a fn (for cpt surfaces) that satisfies:
	\begin{enumerate}[E1)]
		\item if $\Sigma=\Sigma'\sqcup\Sigma''$, then $\chi(\Sigma)=\chi(\Sigma')+\chi(\Sigma'')$
		\item if $\Sigma$ obtained from $\Sigma'$ by gluing two $\p$ circles, then
			$\chi(\Sigma)=\chi(\Sigma')$
		\item $\chi(D^{2})=1$
	\end{enumerate}
\end{defn} \

\begin{xmp}
	since $S^{2}=D^{2}\cup_{\p}D^{2}$, $\chi(S^{2})=2$.
\end{xmp} \

\begin{crll}[title=Hairy Ball Theorem]
	every VF on $S^{2}$ has a zero
\end{crll}
